\lesson{batch}{重复任务先小批试跑，再扩大}
\diagram{p_batch}
多个文件或任务需要相同处理时，先把单项流程做稳定，再扩大规模。批量处理是重复同一套规则；只有各项互不依赖、不会相互覆盖时，才进一步并行。
\topic{从小批结果确定整批规则}
\begin{enumerate}
\item \textbf{先定清单和格式。}列出每项输入、预期输出、保存位置与完成条件，使用同一模板，保留原始材料。
\item \textbf{先处理少量代表项。}包含常见情况和容易失败的情况。检查内容与格式；AI 参与判断时，在测试副本上重复少量代表项，观察结论与质量是否稳定，再决定扩量。
\item \textbf{逐项记录并只补失败。}每项保留成功、失败、跳过或待确认状态及原因。规则、输入或运行环境改变后，需重查受影响项；同类错误反复出现时先停下修正流程。
\end{enumerate}
\begin{prompt}{小批试跑与批量执行}
按[规则和模板]处理[任务清单]，输出到[位置]；允许操作与预算：[范围]。先试少量常见项和高风险项，展示结果与证据；AI参与判断时重复代表项查稳定性。规则或完成条件不清先补齐，缺工具说明，不声称已处理。\par
为每项设唯一标识、输出格式、来源和完成条件。仅并行互不依赖项，避免覆盖原件和重复外部操作；重复试验只用副本或可安全重复的操作。记录状态、结果位置及失败原因；结果未知先核对，不盲目重试。续做未完成项，输入、规则或环境变化时重查受影响项。\par
逐项执行必要检查，再抽查内容；抽查不能替代逐项检查。试跑达标且在授权内再扩量，同类失败重复出现则暂停。汇总完整清单与缺口，比较总耗时、人工检查和返工，不只报告生成速度。
\end{prompt}
\principle{先修正流程，再扩大规模}{先试少量，是为了在错误被重复很多次之前修正方法。看完试跑结果，再确认格式和判断标准；扩大后按清单核对成功项与缺口，失败项保留原因。规则改变时重查受影响项；不要只看到几份好结果，就认定整批都已完成。}
